\section*{الفصل \ref{ch:b2:coordination-complexes} --- الفلزات الانتقالية ومعقدات التناسق}

\begin{solution}{exo:b2:coordination-complexes:1}
\ce{Ti^3+} $d^1$، \ce{V^3+} $d^2$، \ce{Cr^3+} $d^3$، \ce{Mn^2+} $d^5$، \ce{Fe^3+} $d^5$، \ce{Co^2+}
$d^7$، \ce{Ni^2+} $d^8$، \ce{Zn^2+} $d^{10}$.
\end{solution}

\begin{solution}{exo:b2:coordination-complexes:2}
الماء 1، وen 2، والإيثانديوات 2، وedta 6، والثيوسيانات 1؛ والثيوسيانات ثنائية الموقع
(عبر S أو N).
\end{solution}

\begin{solution}{exo:b2:coordination-complexes:3}
كبريتات رباعي أمين النحاس(II)؛ وسداسي سيانيدو حديدات(III) البوتاسيوم؛
وكلوريد رباعي أمين ثنائي كلوريدو الكوبالت(III)؛ ورباعي كربونيل النيكل(0).
\end{solution}

\begin{solution}{exo:b2:coordination-complexes:4}
خطي؛ رباعي الأوجه؛ مربع مستوٍ ($d^8$، السلسلة الثالثة)؛ ثماني الأوجه.
\end{solution}

\begin{solution}{exo:b2:coordination-complexes:5}
اثنان: cis (الكلوريدان متجاوران) وtrans (متقابلان). وفي رباعي الأوجه لن يكون
إلا متماكب واحد.
\end{solution}

\begin{solution}{exo:b2:coordination-complexes:6}
\ce{Cr(CO)6} 18؛ \ce{Fe(CO)5} 18؛ \ce{Ni(CO)4} 18؛ \ce{V(CO)6}: $5 + 12 = 17$؛
\ce{Mn2(CO)10} 18 لكل ذرة منغنيز. ويخالف \ce{V(CO)6} القاعدة: فهو \omterm{def:b2:diatomic-mos:magnetism}{بارامغناطيسي} 
ويُختزل بسهولة إلى \ce{[V(CO)6]-}، الذي له 18.
\end{solution}

\begin{solution}{exo:b2:coordination-complexes:7}
الفيروسين: Fe(II) $d^6$ + 12 = 18. والكوبالتوسين: Co(II) $d^7$ + 12 = 19. ويفقد الكوبالتوسين
إلكترونه الزائد بسهولة، فيعطي أيون الكوبالتوسينيوم \ce{[Co(C5H5)2]+}، الذي له 18.
\end{solution}

\begin{solution}{exo:b2:coordination-complexes:8}
تحل ربيطة edta واحدة محل ست جزيئات ماء: \ce{[Ni(H2O)6]^2+ + edta^4- -> [Ni(edta)]^2- +
6H2O}، سبعة جسيمات من اثنين. \omterm{def:b2:reaction-free-energy:reaction-entropy}{فإنتروبيا التفاعل} موجبة بقوة، 
و$K$ كبير جدًا. أما ست ربائط منفصلة فتحل محل ست جزيئات ماء دون أي زيادة في
عدد الجسيمات.
\end{solution}

\begin{solution}{exo:b2:coordination-complexes:9}
نتريتو-$\kappa N$ (\ce{Co-NO2}، متماكب النيترو، أصفر) ونتريتو-$\kappa O$
(\ce{Co-O-N=O}، متماكب النتريتو، أحمر).
\end{solution}

\begin{solution}{exo:b2:coordination-complexes:10}
ثلاثة: المتماكب trans (لاكيرالي: له مستويات مرآة) والمتماكبان المرآويان للمتماكب
cis (يلتف المخلبان والكلوريدان في اتجاه أو في الآخر).
\end{solution}

\begin{solution}{exo:b2:coordination-complexes:11}
السداسي المستوي: B في الموضعين 1،2 (أورثو)، و1،3 (ميتا)، و1،4 (بارا): ثلاثة متماكبات.
والمنشور الثلاثي: B على حرف من حروف أحد المثلثين، أو على حرف عمودي، أو قطريًا على
وجه مستطيل: ثلاثة. ويعطي ثماني الأوجه اثنين. وعدم إيجاد أكثر من اثنين، رغم
محاولات كثيرة ومركبات عدة، يستبعد الهندستين الأخريين --- ما لم يكن
متماكب ثالث قد فات.
\end{solution}

\begin{solution}{exo:b2:coordination-complexes:12}
\ce{[RhCl(PPh3)3]}: Rh(I) $d^8$ + 4 × 2 = 16. \ce{[IrH(CO)(PPh3)3]}: Ir(I) $d^8$ + 5 × 2 =
18. ومعقد الروديوم، بإلكتروناته الستة عشر، يمكنه استقبال إلكترونين إضافيين (الإضافة التأكسدية،
\cref{ch:b2:catalytic-cycles}).
\end{solution}

\begin{solution}{pb:b2:coordination-complexes:1}
\textbf{1.} الكلوريدات الواقعة خارج \omterm{def:b2:coordination-complexes:coordination-sphere}{كرة التناسق}، الحرة في المحلول.
\textbf{2.} A: 0؛ B: 1؛ C وD: 2.
\textbf{3.} A: $[\ce{Co(NH3)6}]^{3+}$ + 3 \ce{Cl-}: أربعة أيونات؛ B: 3؛ C وD: 2. وهي متوافقة.
\textbf{4.} +3.
\textbf{5.} $d^6$.
\textbf{6.} ست في كل منها: 6 \ce{NH3}؛ 5 \ce{NH3} + 1 Cl؛ 4 \ce{NH3} + 2 Cl.
\textbf{7.} \ce{[Co(NH3)6]Cl3}؛ \ce{[CoCl(NH3)5]Cl2}؛ \ce{[CoCl2(NH3)4]Cl} (مرتين).
\textbf{8.} كلوريد سداسي أمين الكوبالت(III)؛ وكلوريد خماسي أمين كلوريدو الكوبالت(III).
\textbf{9.} رباعي أمين ثنائي كلوريدو الكوبالت(III).
\textbf{10.} متماكبات هندسية (cis--trans).
\textbf{11.} \ce{[CoCl3(NH3)3]}، معقد متعادل: لا أيونات، ولا كلوريد يترسب
فورًا.
\textbf{12.} اثنان: fac وmer.
\textbf{13.} cis: ذرتا Cl عند $90^\circ$؛ trans: عند $180^\circ$.
\textbf{14.} ثلاثة (مواضع أورثو وميتا وبارا على السداسي).
\textbf{15.} ثلاثة.
\textbf{16.} إنه يدعم ثماني الأوجه، الهندسة الوحيدة من الثلاث التي تتنبأ باثنين.
\textbf{17.} لا: فغياب متماكب ثالث دليل سلبي. ولو وُجد متماكب ثالث
لاستبعد ثماني الأوجه.
\textbf{18.} المركب الذي يكوّن المعقد المخلبي: فالإيثانديوات الثنائية المخلب لا تمتد إلا على موضعين
cis.
\textbf{19.} ثلاثة مخالب en حول الكوبالت، يمتد كل منها على موضعين cis.
\textbf{20.} لا هذا ولا ذاك: لا مستوي مرآة ولا مركز تناظر.
\textbf{21.} متماكبان مرآويان، $\Delta$ و$\Lambda$.
\textbf{22.} أن الربائط الست مرتبة في ثلاثة أبعاد، كما في
ثماني الأوجه؛ فالترتيب المستوي يعطي أيونًا لاكيرالي.
\textbf{23.} cis: كيرالي (متماكبان مرآويان). trans: لاكيرالي.
\textbf{24.} بيّن أن النشاط الضوئي خاصية للهندسة الجزيئية عمومًا،
لا للكربون، وأكّد ثماني الأوجه.
\textbf{25.} يتنبأ ثماني الأوجه للأيون \ce{[CoCl2(NH3)4]+} بعدد من المتماكبات قدره $\boldsymbol{2}$؛ ويتنبأ
السداسي المستوي والمنشور الثلاثي بعدد قدره $\boldsymbol{3}$.
\end{solution}
